\begin{titlepage}
\centering
\vspace*{2cm}
{\Large Riemann--Hilbert Approach for the\par}
\medskip
{\Large Hirota--Satsuma Coupled KdV System\par}
\vspace{1cm}
{\huge Hirota--Satsuma 耦合 KdV 系统\par}
\medskip
{\huge 的 Riemann--Hilbert 方法\par}
\vspace{1cm}
{\LARGE 逐句双语与完整推导学习笔记\par}
\vspace{0.6cm}
{\Large 全文逐句双语、校勘与推导\par}
\vfill
{\large 原文：Haibing Zhang、Xianguo Geng、Huan Liu\par}
\medskip
Studies in Applied Mathematics, 155 (2025), e70113\par
\href{https://doi.org/10.1111/sapm.70113}{DOI: 10.1111/sapm.70113}\par
\vspace{1cm}
2026 年 10 月 20 日\par
\end{titlepage}

\chapter*{阅读约定与必要背景}
\addcontentsline{toc}{chapter}{阅读约定与必要背景}
本笔记按论文的原章节顺序编排，保留原文的公式、定义、引理、命题、推论、假设与定理编号。
英文段落按原文顺序转录，紧随中文翻译；共用公式排一次。
英文中的数学符号按本工程统一为直立虚数单位、完整分式和语义化命令，不改数学内容。
EN 标识原文，中文字样标识译文；其余中文的注与证明为补充内容。
原文全部引文编号保留，完整著录列在文末。历史背景与方法动机放在注中解释；数学断言的证明紧随对应内容。

\begin{note}
  \textbf{证明详略与记号。}
  推导尽量沿用原文符号，直接代入原来的函数、矩阵和系数；不为缩短一行公式而连续引入换元或简记。
  对后文按谱区域排列的 $M_n,S_n,T_n$，按原文实际展开的情形给出详细证明。
  原文若只列出 $n=1,8$，就详证这两个，不另补其他六个区域的全套公式和计算。
  其余注明“同理”，交代所依据的区域对应或对称关系即可；若条件或步骤有实质区别，再单独说明。
\end{note}

\begin{note}
  \textbf{编号与校勘。}
  原文公式用原编号，补充推导另用 S2.1、S2.2 等编号。
  英文和中文引理分别显示标题、共享编号；证明在两种陈述及其共用公式之后。
  第 2.1 节 $g_1,g_2$ 的印刷式存在可直接核验的相容性问题：正文先保留原式，随后给出校勘注和证明。
  从修正式引入后起，后续计算中的 $V_1,V,G,V_2$ 均使用明确标出的修正版本。
\end{note}

本文研究的方程组见引言中的式 \eqref{eq:1.1}。
这里 $u=u(x,t)$、$v=v(x,t)$ 是实值函数，下标表示偏导。
例如 $u_{xxx}=\partial_x^3u$，$uv_x$ 表示 $u$ 与 $v$ 的空间偏导之积。
在涉及 Lax 相容性时假设这些函数足够光滑，使出现的混合偏导可以交换；讨论空间衰减时，固定存在时间区间内的 $t$。
在这个固定时刻，Schwartz 衰减的含义是对任意非负整数 $a,b$，
\[
\sup_{x\in\mathbb R}(1+|x|)^a
\bigl(|\partial_x^bu(x,t)|+|\partial_x^bv(x,t)|\bigr)<\infty.
\]
这是正散射实际使用的固定时刻性质。时间演化中在任意有界存在时间区间使用一致估计；原文定义 1.1 对无限时间区间的过强要求，在定义后单独说明。

\begin{note}
  \textbf{怎样衔接范老师教材的前六章。}
  Lax 对（Lax pair）是两个关于辅助函数的线性微分方程。
  原非线性方程出现在这两个线性方程的相容条件中。
  这里辅助函数是四阶矩阵；规范变换（gauge transformation）先将零背景下的常系数部分对角化，
  然后去掉自由指数，才得到适合建立积分方程的形式。
  第 2.1 节只需要矩阵代数、线性 ODE、偏导和复数运算；Fredholm 理论在第 2.6 节按实际需要补充。
\end{note}

\begin{note}
  \textbf{符号速查。}
  $k\in\mathbb C$ 是谱参数，作 $x,t$ 微分时保持不变；$\ii^2=-1$。
  $I_n$ 表示 $n$ 阶单位矩阵；$[X,Y]=XY-YX$ 为交换子。
  上标 $\infty$ 表示位势为零时的背景矩阵；下标 1 表示原坐标中的位势扰动，
  下标 2 表示变换坐标中的位势扰动。
  原文星号 $*$ 表示逐项复共轭，不表示共轭转置；逆转置明确写为 $X^{-\mathsf T}=(X^{-1})^{\mathsf T}$。
\end{note}
\begin{note}
  \textbf{怎样阅读原文中的缺口。}
  原文没有展开的计算在对应的 Proof／证明 中补足。若印刷结论过强或条件不足，则先保留原陈述，再区分已证明的有效版本与原文仍需补充的论证；不为错误命题编造证明。
  纯孤子留数条件是作者另行设定的模型，不能当作已经从一般散射数据推导出的定理。
  定理 3.2 是给定 PDE 解的 RH 表示，定理 3.3 是满足指定 RH 条件后的反向重构；两者均不等同于任意数据的全局适定性。
\end{note}
